\section{Mechanical preparation, collisions, and physical records}\label{sec:mechanics}
\subsection{A mechanical realization class}
For any fixed finite particle number, use a configuration domain $\mathcal Q_{N,a}$ in Euclidean space or a periodic quotient, positive masses, and
\begin{equation}
 H_{N,a}(q,p)=\tfrac12p^TM^{-1}p+V_{N,a}(q),
 \quad M=\diag(m_1I_d,\ldots,m_NI_d).
 \label{eq:natural-H}
\end{equation}
The domain can encode walls and hard-core exclusions; its smooth interior obeys
$\dot q=M^{-1}p$, $\dot p=-\nabla V_{N,a}(q)$. Other symplectic systems fit the experiment construction, but the formulas of this section concern \eqref{eq:natural-H}. Existence of an arbitrary-particle or low-density limit is not part of the definition.

\begin{lemma}[Local flow and parameter variation]\label{lem:flow}
Suppose $f_a(z)$ is $C^1$ on a common compact neighborhood of the relevant trajectories, with $\norm{D_zf_a}\le L$ and $\norm{\partial_af_a}\le M_1$. For equal initial states,
\[
 |z_a(t)-z_b(t)|\le M_1|a-b|(e^{Lt}-1)/L.
\]
The quotient is $t$ when $L=0$. The parameter derivative satisfies
\[
 \dot z_1=D_zf_a(z)z_1+\partial_af_a(z),
 \quad z_1(0)=\partial_az_a(0).
\]
If $f$ and the initial state are $C^r$, the flow is $C^r$ on the common interval, with the explicit jet recursion of Appendix~\ref{app:jets}.
\end{lemma}
\begin{proof}
On a short interval with $Lh<1$, the integral map
$z\mapsto z_0+\int_0^tf_a(z(s))ds$ is a contraction on a closed ball staying in the neighborhood. Finite continuation gives the common interval. Subtract the integral equations and apply Gronwall to obtain the displayed bound. Subtract the candidate linearized integral equation from the difference-quotient equation. Uniform continuity of first derivatives and the same bound make its forcing error tend uniformly to zero; Gronwall gives strong convergence of the difference quotients. Repeating this argument, each higher variation solves a linear equation with a polynomial forcing involving lower variations and derivatives of $f$. The polynomial and its norm recursion are given in Appendix~\ref{app:jets}.
\end{proof}

\subsection{Mass-metric reflection}
For a fixed wall $g_a(q)=0$, write $n=\nabla g_a(q)\ne0$. Elastic reflection is
\begin{equation}
 p^+=p^--2\frac{n^TM^{-1}p^-}{n^TM^{-1}n}n.
 \label{eq:reflection}
\end{equation}
The normal may be scaled or reversed without changing the formula.

\begin{proposition}[Energy, reversibility, and boundary flux]\label{prop:reflection}
Reflection \eqref{eq:reflection} is an involution, reverses the mass-normal velocity, and preserves $p^TM^{-1}p/2$. At a regular impact it preserves the absolute normal flux measure together with tangential coordinates. For two hard particles it gives the usual opposite impulses and preserves total momentum and kinetic energy.
\end{proposition}
\begin{proof}
Let $c=n^TM^{-1}p^-$ and $d=n^TM^{-1}n$. Then $n^TM^{-1}p^+=c-2c=-c$. A second application restores $p^-$. Expanding the kinetic quadratic form gives
$p^{+T}M^{-1}p^+=p^{-T}M^{-1}p^- -4c^2/d+4c^2/d$.
In coordinates $v=M^{-1/2}p$, reflection is Euclidean reflection across the plane perpendicular to $M^{-1/2}n$, so its absolute velocity Jacobian is one. The magnitude of normal velocity is unchanged; therefore the incident flux factor is unchanged. Interior Hamiltonian flow preserves volume because its divergence is zero. For a two-particle contact $n$ has blocks $\widehat n,-\widehat n$ and zero elsewhere; the impulse blocks sum to zero. Explicitly,
\[
 v_i^+=v_i^- -\frac{2m_j}{m_i+m_j}[(v_i^--v_j^-)\cdot\widehat n]\widehat n,
\]
with the opposite mass-weighted increment for $j$. The preceding quadratic calculation applies.
\end{proof}

A moving wall in physical time has a different energy balance: relative normal velocities involve wall velocity, and the wall can do work. Parameter differentiation of a family of fixed tables is not the same experiment as physically moving a wall during a trajectory.

\subsection{Preparations and complete records}
A finite nonzero Liouville, energy-shell, or other explicitly specified measure can be normalized to a preparation. A non-equilibrium density is equally allowed. A conserved quantity has to be fixed by a geometric measure or disintegration when its level set has zero ambient volume. Different preparations on the same Hamiltonian generally give different path statistics.

For a nonexplosive finite interval, retain $z_0$, the flight pieces, all event times, pre- and post-impact states, collision labels, and the final state. A cemetery or stopped record is retained when the model has a declared exit. One must prove that excluded events have zero preparation probability before discarding them. For finitely many events, the record lies in a countable union of finite-dimensional Borel products. For countably many events, a suitable sequence space or specified cadlag path space is used; no compactness follows merely from this encoding.

\begin{proposition}[Exact flight--collision balance]\label{prop:balance}
For a piecewise $C^1$ trajectory with finitely many impacts and a $C^1$ test $\varphi(t,z)$ on each flight,
\begin{align}
 \varphi(T,z_T)-\varphi(0,z_0)
 &=\int_0^T(\partial_t\varphi+\{\varphi,H\})(t,z_t)\dd t\notag\\
 &\quad+\sum_{t_k\le T}[\varphi(t_k,z_{t_k}^+)-\varphi(t_k,z_{t_k}^-)].
 \label{eq:path-balance}
\end{align}
\end{proposition}
\begin{proof}
Integrate the ordinary chain rule separately on every open flight interval. Add and subtract the two one-sided values at every impact. All adjacent endpoints telescope, leaving exactly the displayed jump terms. At an impact at $T$, use the declared right-continuous terminal convention. For an infinite-event record the same conclusion requires absolute summability of jumps and convergence of endpoint traces; these are not supplied by the finite calculation.
\end{proof}

The resulting occupation and collision measures satisfy a distributional balance. The converse is not automatic: a measure satisfying a formal balance need not be a realizable mechanical history. Recovery of such measures is a separate problem for later large-deviation papers.

\subsection{A marked stopping clock}
For history updates $h_k=U(h_{k-1},m_k)$, positive holding lengths $R(m_k)$, scale $N$, and $s_k=N^{-1}\sum_{j\le k}R(m_j)$, define
\[
 J_N=\sum_{k=1}^{\ell_N}\delta_{(s_k,h_{k-1},m_k)},\quad Q_N=J_N/N,
\]
\[
 L_N(ds,dh)=\sum_{k=1}^{\ell_N}\one_{[s_{k-1},s_k)}(s)\dd s\,\delta_{h_{k-1}}(dh).
\]
Here $\ell_N$ is a specified stopping count, not an occupation measure. For
$\langle\operatorname{div}_UJ,\phi\rangle=\int[\phi(s,h)-\phi(s,U(h,m))]dJ$,
\begin{equation}
 \partial_sL_N+\operatorname{div}_UJ_N
 =\delta_{(0,h_0)}-\delta_{(s_{\ell_N},h_{\ell_N})}.
 \label{eq:clock-balance}
\end{equation}
Indeed each edge difference is the sum of the time increment at its tail and the jump from tail to head at its right endpoint; summation telescopes. The coefficient of $\operatorname{div}_UQ_N$ is therefore $N$.

For a bounded time-Lipschitz $g$,
\[
 \left|\int g\dd L_N-\int R(m)g(s,h)\dd Q_N\right|
 \le\tfrac12\Lip_s(g)\sum_k(R(m_k)/N)^2.
\]
This follows by integrating $g(s_k-r,h)-g(s_k,h)$ over each holding interval. A macroscopic holding interval is not a weighted point atom. This distinction is retained when returning from collision count to physical time.
\par
